\uuid{gkog}
\exo7id{1855}
\titre{exo7 1855}
\auteur{maillot}
\organisation{exo7}
\datecreate{2001-09-01}
\isIndication{false}
\isCorrection{false}
\chapitre{Equation aux dérivées partielles}
\sousChapitre{Equation aux dérivées partielles}
\cpge{mpsi-mp2i.fonctions-deux-variables.b, pcsi.fonctions-deux-variables.b, ptsi.fonctions-deux-variables.b, mpsi-mp2i.fonctions-deux-variables.c, pcsi.fonctions-deux-variables.c, ptsi.fonctions-deux-variables.c, bcpst1.fonctions-deux-variables.c, mp-mpi.calcul-differentiel-optimisation.a}
\module{Equation différentielle}
\niveau{L3}
\difficulte{}

\contenu{
\texte{
Soit $D=\{(x,y)\in \R^2,\, x>0\}$. On cherche les fonctions $f \in {\cal
C}^1(D,\R)$ qui v\'erifient
$$(E) \qquad x \frac{\partial f}{\partial x} + y  \frac{\partial
f}{\partial y} =0 \quad \forall (x,y) \in D.$$
}
\begin{enumerate}
    \item \question{V\'erifier que $\varphi(x,y)=y/x$ est solution de (E).}
    \item \question{Soit $g \in {\cal C}^1(\R, \R)$. Montrer que $g \circ \varphi$ est
   solution de $(E)$.}
    \item \question{Soit $f$ une solution de $(E)$. Montrer que $f(u, uv)$ ne d\'epend
   que de $v$.}
    \item \question{Donner l'ensemble des solutions de $(E)$.}
\end{enumerate}
}
